\chapter{Chronic Pain}

\section*{Definition}
Pain persisting or recurring for longer than three months, beyond typical tissue healing time; classified as primary (no identifiable underlying cause, e.g., fibromyalgia) or secondary (due to an underlying condition such as arthritis or neuropathy).

\section*{What Happens in Your Body}
Chronic pain involves maladaptive neuroplasticity and central sensitization -- persistent nociceptive input leads to altered processing in the dorsal horn, thalamus, and cortex with reduced descending inhibitory control. Prolonged glial activation releases pro-inflammatory immune signaling proteins (IL-6, TNF-alpha) that maintain hyperexcitability, and structural changes in gray matter volume correlate with pain chronicity.

\section*{Causes}
\begin{itemize}
  \item Nociceptive: Osteoarthritis, rheumatoid arthritis, mechanical low back pain, chronic post-surgical pain
  \item Neuropathic: Diabetic neuropathy, postherpetic neuralgia, trigeminal neuralgia, spinal stenosis, complex regional pain syndrome
  \item Nociplastic: Fibromyalgia, chronic tension-type headache, irritable bowel syndrome, interstitial cystitis, temporomandibular disorder
  \item Mixed: Cancer pain (tumor infiltration + nerve compression), chronic migraine, chronic pelvic pain
\end{itemize}

\section*{When to See a Doctor}
See your doctor if:
\begin{itemize}
  \item Pain interferes with sleep, work, mobility, or basic activities of daily living
  \item Red flags: Unintentional weight loss, nocturnal pain, fever, history of malignancy, recent trauma, progressive neurologic deficits
  \item Pain accompanied by mood disturbance, suicidal ideation, or substance use concerns
  \item Inadequate response to initial treatment measures after 4-6 weeks
\end{itemize}

\section*{Related Symptoms}
\begin{itemize}
  \item Fatigue
  \item Depression and anxiety
  \item Sleep disturbance
  \item Cognitive dysfunction (brain fog)
  \item Reduced physical function
\end{itemize}
\textbf{Important:} Chronic pain is best managed through a biopsychosocial model combining pharmacologic, physical, and psychological interventions.

\section*{Self-Care / Home Management}
\begin{itemize}
  \item Maintain regular physical activity within pain limits: walking, swimming, stretching, or graded exercise programs.
  \item Practice paced activity (activity-rest cycling) to avoid flare-ups from overexertion.
  \item Use heat or cold therapy, meditation, mindfulness, and deep breathing for flare-up management.
  \item Establish consistent sleep hygiene and consider cognitive behavioral therapy (CBT) for pain coping strategies.
\end{itemize}

\section*{How Your Doctor Will Evaluate You}
\begin{itemize}
  \item Detailed pain history: onset, location, quality, intensity (0-10 scale), timing, aggravating/alleviating factors.
  \item Pain mapping and physical exam assessing for allodynia, hyperalgesia, nerve tension signs.
  \item Imaging (X-ray, MRI) and electrodiagnostic studies (EMG/NCS) guided by clinical suspicion.
  \item Psychosocial assessment for mood disorders, catastrophizing, kinesiophobia, and functional goals.
\end{itemize}

\section*{Treatment Options}
\begin{itemize}
  \item First-line pharmacotherapy: NSAIDs (nociceptive), gabapentinoids or SNRIs (neuropathic), tricyclic antidepressants (central sensitization).
  \item Physical therapy, occupational therapy, and graded motor imagery for functional restoration.
  \item Interventional: Epidural steroid injections, nerve blocks, radiofrequency ablation, spinal cord stimulation.
  \item Psychological: CBT, acceptance and commitment therapy (ACT), biofeedback, multidisciplinary pain rehabilitation.
\end{itemize}

\section*{References}
\begin{thebibliography}{99}
\bibitem{chronic_pain:ref1} Dowell D, Ragan KR, Jones CM, et al. ``CDC clinical practice guideline for prescribing opioids for pain.'' MMWR Recomm Rep. 2022;71(3):1-95.
\bibitem{chronic_pain:ref2} Cohen SP, Vase L, Hooten WM. ``Chronic pain: an update on burden, best practices, and new advances.'' Lancet. 2021;397(10289):2082-2097.
\bibitem{chronic_pain:ref3} Woolf CJ. ``Central sensitization: implications for the diagnosis and treatment of pain.'' Pain. 2011;152(3 Suppl):S2-S15.
\bibitem{chronic_pain:ref4} Treede RD, Rief W, Barke A, et al. ``Chronic pain as a symptom or a disease: the IASP classification.'' Pain. 2019;160(1):19-27.
\bibitem{chronic_pain:ref5} Benzon HT, Raja SN, Fishman SM, et al. ``Essentials of Pain Medicine.'' 4th ed. Elsevier; 2018.
\bibitem{chronic_pain:ref6} Rosenquist RW. ``Evaluation of chronic pain in adults.'' UpToDate. Wolters Kluwer; 2025.
\end{thebibliography}
